\documentclass[preview]{standalone}

\usepackage{amsmath}
\usepackage{amssymb}
\usepackage{stellar}
\usepackage{definitions}

\begin{document}

\id{ring-substructures-exercises}
\genpage

\section{Exercises}

\begin{snippetexercise}{upper-triangular-matrix-ring-ideals-exercise}{Subrings and ideals in a matrix ring}
    Consider the \subring
    \[
        A=
        \left\{
            \begin{pmatrix}a&c\\0&b\end{pmatrix}
            \suchthat a,b,c\in\integers
        \right\}
        \subseteq\matrices_2(\integers).
    \]
    First verify that \(A\) is a ring. Then determine which of the
    following subsets are two-sided \ideal[ideals] of \(A\):
    \[
        S_1=
        \left\{
            \begin{pmatrix}a&0\\0&a\end{pmatrix}
            \suchthat a\in\integers
        \right\},
    \]
    \[
        S_2=
        \left\{
            \begin{pmatrix}a&c\\0&0\end{pmatrix}
            \suchthat a,c\in\integers
        \right\},
        \qquad
        S_3=
        \left\{
            \begin{pmatrix}2h&4\ell\\0&2k\end{pmatrix}
            \suchthat h,k,\ell\in\integers
        \right\}.
    \]
\end{snippetexercise}

\begin{snippetsolution}{upper-triangular-matrix-ring-ideals-solution}{Subrings and ideals in a matrix ring}
    The zero and identity matrices belong to \(A\). If
    \[
        M=\begin{pmatrix}a&c\\0&b\end{pmatrix},
        \qquad
        N=\begin{pmatrix}a'&c'\\0&b'\end{pmatrix},
    \]
    then
    \[
        M-N=
        \begin{pmatrix}a-a'&c-c'\\0&b-b'\end{pmatrix}
        \in A
    \]
    and
    \[
        MN=
        \begin{pmatrix}aa'&ac'+cb'\\0&bb'\end{pmatrix}
        \in A.
    \]
    Thus \(A\subringleq\matrices_2(\integers)\).

    The subset \(S_1\) contains \(I_2=1_A\) but is not all of \(A\);
    hence it cannot be an ideal.

    The subset \(S_2\) is an additive subgroup. For \(M\in A\) and
    \(N\in S_2\),
    \[
        MN=
        \begin{pmatrix}aa'&ac'\\0&0\end{pmatrix}\in S_2,
        \qquad
        NM=
        \begin{pmatrix}a'a&a'c+c'b\\0&0\end{pmatrix}\in S_2.
    \]
    Therefore \(S_2\idealin A\).

    Finally,
    \[
        \begin{pmatrix}2&4\\0&2\end{pmatrix}\in S_3,
        \qquad
        \begin{pmatrix}0&1\\0&0\end{pmatrix}\in A,
    \]
    but their product is
    \[
        \begin{pmatrix}0&2\\0&0\end{pmatrix}\notin S_3.
    \]
    Thus \(S_3\) is not a two-sided ideal.
\end{snippetsolution}

\begin{snippetexercise}{two-point-integer-function-ring-ideals-exercise}{Ideals in a two-point function ring}
    Let \(X=\{x,y\}\), and let
    \[
        A=\{f\colon X\fromto\integers\}
    \]
    have pointwise addition and multiplication. Show that \(A\) is a
    commutative ring, and determine whether the following subsets are
    ideals:
    \[
        S_1=\{f\in A\suchthat
            2\divides f(x)\text{ or }2\divides f(y)\},
    \]
    \[
        S_2=\{f\in A\suchthat 2f(x)=3f(y)+4\},
        \qquad
        S_3=\{f\in A\suchthat4f(x)-2f(y)=0\}.
    \]
\end{snippetexercise}

\begin{snippetsolution}{two-point-integer-function-ring-ideals-solution}{Ideals in a two-point function ring}
    The ring axioms follow pointwise from those of \(\integers\). The
    zero and unit are the constant functions \(0\) and \(1\), and
    commutativity is inherited pointwise.

    Define \(f,g\in A\) by
    \[
        (f(x),f(y))=(1,0),
        \qquad
        (g(x),g(y))=(0,1).
    \]
    Both lie in \(S_1\), but \(f+g\) has values \((1,1)\) and does not
    lie in \(S_1\). Hence \(S_1\) is not an additive subgroup.

    The zero function does not lie in \(S_2\), because \(0\neq4\).
    Thus \(S_2\) is not an ideal.

    The equation defining \(S_3\) is homogeneous, so \(S_3\) is an
    additive subgroup. Nevertheless, take
    \[
        (f(x),f(y))=(1,2),
        \qquad
        (g(x),g(y))=(1,0).
    \]
    Then \(f\in S_3\), but
    \[
        4(fg)(x)-2(fg)(y)=4\neq0,
    \]
    so \(fg\notin S_3\). Therefore none of \(S_1,S_2,S_3\) is an
    ideal.
\end{snippetsolution}

\begin{snippetexercise}{function-ring-vanishing-and-ideal-values-exercise}{Vanishing and ideal-valued functions}
    Let \(A\) be a commutative \ring, let \(I\idealin A\) be proper,
    let \(X\) be a nonempty \set, and let \(Y\subseteq X\). In the ring
    \[
        B=A^X=\{f\colon X\fromto A\}
    \]
    with pointwise operations, determine which of the following are
    ideals:
    \[
        S_1=\{f\in B\suchthat f(y)=0_A
            \text{ for every }y\in Y\},
    \]
    \[
        S_2=\{f\in B\suchthat f(y)=1_A
            \text{ for every }y\in Y\},
    \]
    \[
        S_3=\{f\in B\suchthat f(y)=c
            \text{ for every }y\in Y\},
        \qquad c\in I\text{ fixed},
    \]
    \[
        S_4=\{f\in B\suchthat f(y)\in I
            \text{ for every }y\in Y\}.
    \]
\end{snippetexercise}

\begin{snippetsolution}{function-ring-vanishing-and-ideal-values-solution}{Vanishing and ideal-valued functions}
    If \(Y=\emptyset\), all four conditions are vacuous, so
    \[
        S_1=S_2=S_3=S_4=B;
    \]
    all four are ideals.

    Suppose \(Y\neq\emptyset\). If \(f,g\in S_1\) and \(h\in B\), then
    on \(Y\)
    \[
        (f-g)(y)=0_A,
        \qquad
        (hf)(y)=h(y)0_A=0_A.
    \]
    Hence \(S_1\idealin B\).

    The zero function does not belong to \(S_2\), so \(S_2\) is not an
    ideal. Similarly, \(S_3=S_1\) when \(c=0_A\); if \(c\neq0_A\), the
    zero function does not belong to \(S_3\), so it is not an ideal.

    Finally, if \(f,g\in S_4\) and \(h\in B\), then for every \(y\in Y\)
    \[
        f(y)-g(y)\in I,
        \qquad
        h(y)f(y)\in I.
    \]
    Thus \(S_4\idealin B\). It is proper because \(I\) is proper,
    \(Y\neq\emptyset\), and the constant function \(1_A\) does not
    belong to \(S_4\).
\end{snippetsolution}

\end{document}
